% GENERATED FILE. Edit canonical lessons, then run npm run generate:books.
% canonical-source-hash: fnv1a64:8a9c544dbf61f4c4
% canonical-lessons: IT-C334-soffriggere, IT-C334-rosolare, IT-C334-friggere, IT-C334-sfumare, IT-C334-soffritto

\chapter{In padella}
\label{ch:it-in-padella}
\hlchaptermodality{334}

\begin{chapteropening}
\textbf{By the end of this chapter:} \emph{I can describe frying: the gentle fry of the onion, browning, frying and cooking off the wine.}

Complete the last lesson of chapter 334: I can describe frying: the gentle fry of the onion, browning, frying and cooking off the wine.
\end{chapteropening}

\section[soffriggere]{soffriggere — to fry gently, to sauté}

\label{lesson:IT-C334-soffriggere}



\begin{quote}
Before the new one: say the Italian for the lid, then the Italian for the shallow pan.
\end{quote}

\subsection*{You'll want to know: soffriggere}
\textbf{soffriggere} — \textquotedblleft{}to fry gently, to sauté\textquotedblright{}.

\begin{quote}
Soffriggi la cipolla nell'olio.
\end{quote}

\begin{quote}
\emph{Fry the onion gently in the oil.}
\end{quote}

Low heat, a few minutes, a little oil.

\begin{cousinweb}
\emph{soffriggere} joins \emph{so-} \textquotedblleft{}under\textquotedblright{}, here \textquotedblleft{}gently\textquotedblright{}, from Latin \emph{sub}, and \emph{friggere} \textquotedblleft{}to fry\textquotedblright{}, from Latin \emph{frīgere} (EtymDB).
\end{cousinweb}

\begin{grammarlens}[title={What you've built}]
One more word for telling, explaining and describing at B1.
\end{grammarlens}

\subsection*{Your turn}
\begin{itemize}
\raggedright
  \item \emph{Say it:} \emph{soffriggere}
  \item \emph{Say it:} \emph{soffriggere}, once more
  \item \emph{Say it:} a sentence of your own with \emph{soffriggere}
  \item \emph{Recall:} say the Italian for the lid, then the Italian for the shallow pan, then say \emph{soffriggere} again
\end{itemize}

\begin{tcolorbox}[breakable,colback=teal!4,colframe=teal!35!black,title={Before you move on}]
What does soffriggere mean? (\textquotedblleft{}To fry gently, to sauté\textquotedblright{}.) Say it once more.
\end{tcolorbox}
\section[rosolare]{rosolare — to brown}

\label{lesson:IT-C334-rosolare}



\begin{quote}
Before the new one: say the Italian for the shallow pan, then the Italian for to fry gently.
\end{quote}

\subsection*{You'll want to know: rosolare}
\textbf{rosolare} — \textquotedblleft{}to brown\textquotedblright{}.

\begin{quote}
Rosola la carne da tutti e due i lati.
\end{quote}

\begin{quote}
\emph{Brown the meat on both sides.}
\end{quote}

Until it turns a golden colour.

\begin{cousinweb}
The origin of \emph{rosolare} is not settled; it is usually linked with \emph{rosso} \textquotedblleft{}red\textquotedblright{}, from Latin \emph{russus} (Lewis \& Short).
\end{cousinweb}

\begin{grammarlens}[title={What you've built}]
One more word for telling, explaining and describing at B1.
\end{grammarlens}

\subsection*{Your turn}
\begin{itemize}
\raggedright
  \item \emph{Say it:} \emph{rosolare}
  \item \emph{Say it:} \emph{rosolare}, once more
  \item \emph{Say it:} a sentence of your own with \emph{rosolare}
  \item \emph{Recall:} say the Italian for the shallow pan, then the Italian for to fry gently, then say \emph{rosolare} again
\end{itemize}

\begin{tcolorbox}[breakable,colback=teal!4,colframe=teal!35!black,title={Before you move on}]
What does rosolare mean? (\textquotedblleft{}To brown\textquotedblright{}.) Say it once more.
\end{tcolorbox}
\section[friggere]{friggere — to fry}

\label{lesson:IT-C334-friggere}



\begin{quote}
Before the new one: say the Italian for to fry gently, then the Italian for to brown.
\end{quote}

\subsection*{You'll want to know: friggere}
\textbf{friggere} — \textquotedblleft{}to fry\textquotedblright{}.

\begin{quote}
Ho fritto il pesce in tanto olio.
\end{quote}

\begin{quote}
\emph{I fried the fish in plenty of oil.}
\end{quote}

Past participle \emph{fritto}; \emph{il fritto misto} is a plate of mixed fried food.

\begin{cousinweb}
\emph{friggere} is inherited from Latin \emph{frīgere} \textquotedblleft{}to roast, to fry\textquotedblright{} (EtymDB).
\end{cousinweb}

\begin{grammarlens}[title={What you've built}]
One more word for telling, explaining and describing at B1.
\end{grammarlens}

\subsection*{Your turn}
\begin{itemize}
\raggedright
  \item \emph{Say it:} \emph{friggere}
  \item \emph{Say it:} \emph{friggere}, once more
  \item \emph{Say it:} a sentence of your own with \emph{friggere}
  \item \emph{Recall:} say the Italian for to fry gently, then the Italian for to brown, then say \emph{friggere} again
\end{itemize}

\begin{tcolorbox}[breakable,colback=teal!4,colframe=teal!35!black,title={Before you move on}]
What does friggere mean? (\textquotedblleft{}To fry\textquotedblright{}.) Say it once more.
\end{tcolorbox}
\section[sfumare]{sfumare — to deglaze, to let the wine cook off; to fade}

\label{lesson:IT-C334-sfumare}



\begin{quote}
Before the new one: say the Italian for to brown, then the Italian for to fry.
\end{quote}

\subsection*{You'll want to know: sfumare}
\textbf{sfumare} — \textquotedblleft{}to deglaze, to let the wine cook off; to fade\textquotedblright{}.

\begin{quote}
Sfuma con un bicchiere di vino bianco.
\end{quote}

\begin{quote}
\emph{Deglaze with a glass of white wine.}
\end{quote}

A risotto word; colours and sounds \emph{sfumano} too, they fade.

\begin{cousinweb}
\emph{sfumare} is built on \emph{fumo} \textquotedblleft{}smoke\textquotedblright{}, from Latin \emph{fūmus} (Lewis \& Short): to go off like smoke.
\end{cousinweb}

\begin{grammarlens}[title={What you've built}]
One more word for telling, explaining and describing at B1.
\end{grammarlens}

\subsection*{Your turn}
\begin{itemize}
\raggedright
  \item \emph{Say it:} \emph{sfumare}
  \item \emph{Say it:} \emph{sfumare}, once more
  \item \emph{Say it:} a sentence of your own with \emph{sfumare}
  \item \emph{Recall:} say the Italian for to brown, then the Italian for to fry, then say \emph{sfumare} again
\end{itemize}

\begin{tcolorbox}[breakable,colback=teal!4,colframe=teal!35!black,title={Before you move on}]
What does sfumare mean? (\textquotedblleft{}To deglaze, to let the wine cook off; to fade\textquotedblright{}.) Say it once more.
\end{tcolorbox}
\section[il soffritto]{il soffritto — the soffritto, the gently fried base of onion and vegetables}

\label{lesson:IT-C334-soffritto}



\begin{quote}
Before the new one: say the Italian for to fry, then the Italian for to deglaze.
\end{quote}

\subsection*{You'll want to know: il soffritto}
\textbf{il soffritto} — \textquotedblleft{}the soffritto, the gently fried base of onion and vegetables\textquotedblright{}.

\begin{quote}
Il soffritto è la base di molte ricette.
\end{quote}

\begin{quote}
\emph{Soffritto is the base of many recipes.}
\end{quote}

Onion, carrot and celery, chopped and fried slowly.

\begin{cousinweb}
\emph{soffritto} is the past participle of \emph{soffriggere} \textquotedblleft{}to fry gently\textquotedblright{}, from Latin \emph{frīgere} \textquotedblleft{}to fry\textquotedblright{} (EtymDB).
\end{cousinweb}

\begin{grammarlens}[title={What you've built}]
One more word for telling, explaining and describing at B1.
\end{grammarlens}

\subsection*{Your turn}
\begin{itemize}
\raggedright
  \item \emph{Say it:} \emph{il soffritto}
  \item \emph{Say it:} \emph{il soffritto}, once more
  \item \emph{Say it:} a sentence of your own with \emph{il soffritto}
  \item \emph{Recall:} say the Italian for to fry, then the Italian for to deglaze, then say \emph{il soffritto} again
\end{itemize}

\begin{tcolorbox}[breakable,colback=teal!4,colframe=teal!35!black,title={Before you move on}]
What does il soffritto mean? (\textquotedblleft{}The soffritto, the gently fried base of onion and vegetables\textquotedblright{}.) Say it once more.
\end{tcolorbox}
